% Page 106

\seite{106}

Jakob Goebel, Sohn des Peter Goebel. Leider folgte einem\ob
vor zwei Jahren erlebten Blitzunglücksfall, bei welchem\ob
seine beiden Pferde tot blieben — seit unferem\ob
Monaten an der Beschlagung darauf. Nachdem er bereits\ob
6 Wochen in einem hiesigen Gasthofe untergebracht und\ob
behandelt worden längerer Zeit zu Hause war, um\ob
die er weiter wegen Gefährdung der öffentlichen Sicher\ohb
heit auf Einverständnis der Einweisung in eine Zwangs\ohb
weise durch Polizei und Gendarm ins Klinik nach Trier.\ob
\unsicher gebracht.

10.\ August  Bei dem gestrigen Volksentscheid, Landtagsauflö\ohb
sung wurden hier 12 Ja-Stimmen, 9 Nein-Stimmen\ob
und 4 ungültige Stimmen bei 137 Wahlberechtigten\ob
abgegeben worden. Volksbegehrens-Einbringung: Rothe Gerhard.

20.\ August  Da infolge der anhaltenden heißen Witterung\ob
die Getreideernte gefährdet erscheint, gingen heute\ob
7 Pflanzen unter Führung der Pfarrer Seffern\ob
prozessionsweise nach Echtlung\unsicher{} in der dor\ohb
tigen Wallfahrtskirche um besseres Wetter zu beten.

17.\ September Mit dem heutigen Tage beginnen die Vierwö\ohb
chigen Herbstferien. Die Schulen beginnen pflichtmäßig\ob
21.\ Oktober.

1.\ Oktober Vor einiger Zeit wurde dem Landwirte Jakob\ob
Hoffmann von hier ein ganz geringfügiger aus dem\ob
Felde 2 Pfund junge Eichen abgeschnitten und gegen\ob
ihn Seitens der Polizei Anzeige wegen\unsicher{} Entpflegen\ob
von Bäume erhoben\unsicher. Hierüber hat sich derart aufge\ohb
stellt, daß der Lehrer einen\unsicher{} stellt wie er in\ob
letzten Jahren anhalten worden ist nicht war, daß sein\ob
Sohn Wolfshorn am 20.\unsicher geschädigt\unsicher{} wurde,\ob
so war in fremden Hause\unsicher, war also\unsicher{} Tage\ob
\unsicher Gefängnis und Entlassung am 2\unsicher{} hiesigen\ob
Vollzug einbrachte.

\pend
\abschnitt{1. Dezember Lehrerstelle Mahrbach in Ließburg [?] dem 1. November 1931}
\pstart

\pend
